% Generated draft: requires mathematical and source review.
\documentclass[11pt]{article}
\usepackage[margin=25mm]{geometry}
\usepackage{fontspec}
\setmainfont{Latin Modern Roman}
\usepackage{amsmath,amssymb,mathtools,unicode-math}
\setmathfont{Latin Modern Math}
\usepackage{enumitem,needspace,xurl,hyperref}
\setlength{\parindent}{0pt}
\setlength{\parskip}{5pt}
\setlength{\emergencystretch}{3em}
\newcommand{\N}{\mathbb{N}}
\newcommand{\Z}{\mathbb{Z}}
\newcommand{\Q}{\mathbb{Q}}
\newcommand{\R}{\mathbb{R}}
\newcommand{\C}{\mathbb{C}}
\newcommand{\F}{\mathbb{F}}
\newcommand{\A}{\mathbb{A}}
\newcommand{\PP}{\mathbb{P}}
\newcommand{\E}{\mathbb{E}}
\newcommand{\eps}{\varepsilon}
\newcommand{\dd}{\,\mathrm{d}}
\begin{document}
\section*{UnsolvedMath ID: 4600034}
\textbf{Source problem code:} catalogue entry
\subsection{Definitions and mathematical statement}
A one-sided shift of finite type is a space of right-infinite paths in a finite directed graph, with the shift deleting the first edge; irreducibility means that the presenting graph may be taken strongly connected. A cover is a continuous surjective shift-commuting map \(\pi:Y\to X\). Its quotient relation is
\[
R_\pi=\{(u,v)\in Y^2:\pi(u)=\pi(v)\}.
\]
For a labeled-graph presentation, left-resolving means that distinct edges entering the same vertex have distinct labels. The extracted record does not specify the construction making the covers canonical.

If the two quotient relations are presented on a common cover shift \((Y,\sigma)\), the proposed equivalence is the existence of
\[
h\in\operatorname{Aut}(Y,\sigma),\qquad (h\times h)(R_{\pi_1})=R_{\pi_2},
\]
where \(h\) is a homeomorphism satisfying \(h\sigma=\sigma h\). Equivalently, for all \(u,v\in Y\),
\[
\pi_1(u)=\pi_1(v)\quad\Longleftrightarrow\quad
\pi_2(h(u))=\pi_2(h(v)).
\]
The question is to give a terminating decision algorithm which, from the intended finite presentations of the two canonical left-resolving irreducible covers of a one-sided irreducible SFT, returns yes exactly when such an automorphism exists. This is a decision problem, not just enumeration of automorphisms or a semidecision procedure.

If the covers have distinct domains \(Y_1,Y_2\), an isomorphism between domains or a specified identification with a common domain is needed before the word automorphism has this meaning. The source supplies neither that convention nor a finite encoding of the canonical covers and their quotient relations. Its stated base is an SFT; it is not silently changed to a sofic shift. The supplied research summary describes the decision question as open, mentioning uniqueness of a canonical cover up to a quotient-respecting right isomorphism. That statement does not itself decide the requested equivalence.
\subsection{Short English statement}
Find an algorithm deciding whether an automorphism carries one canonical cover's quotient relation to the other's. The canonical construction, common domain, and finite input presentation must first be fixed.
\subsection{Sources}
\begin{itemize}
\item Local UnsolvedMath registry, permanent ID \(4600034\), exact source code \texttt{catalogue entry}, ``Equivalence of canonical covers.''
\item Boyle, \emph{Open Problems in Symbolic Dynamics} (2008), Problem/Question/Conjecture 21.2; \url{https://www.math.umd.edu/~mboyle/papers/openfinalsub3nov2007.pdf}. Recorded access 2026-07-29; pdf-text extraction.
\item Supplied OPEN-TRIAGE research summary relates the problem to canonical-cover rigidity and decidability of sofic conjugacy, but supplies no precise additional theorem or publication citation.
\end{itemize}
\end{document}
